\documentclass[10pt,a4paper]{amsart}
\usepackage[margin=19mm]{geometry}
\usepackage{amsmath,amssymb,mathtools}
\usepackage[scr=rsfs,cal=euler]{mathalfa}
\usepackage{xfrac}
\usepackage[unicode,hidelinks,bookmarks=false]{hyperref}
\newcommand{\ie}{i.e.\ }
\newcommand{\cf}{cf.\ }
\newcommand{\resp}{respectively\ }
\newcommand{\Subsection}{\subsection}
\providecommand{\nobreakdash}{\nobreak\hbox{-}}
\setlength{\emergencystretch}{2em}
\multlinegap0pt
\usepackage{bm}
\usepackage{amssymb}
\usepackage{float}
\usepackage{enumerate}
\usepackage{algorithm}\usepackage{algpseudocode}
\newtheorem{proposition}{Proposition}
\newcommand{\cO}{\mathcal{O}}
\title{From molecular dynamics to kinetic models: data-driven generalized collision operators in 1D3V plasmas}
\author{Yue Zhao, Guosheng Fu, Huan Lei}
\date{Source: arXiv:2603.27828v1 (CC BY 4.0)}
\begin{document}
\maketitle

\noindent\emph{Source and licence.} From molecular dynamics to kinetic models: data-driven generalized collision operators in 1D3V plasmas. Version arXiv:2603.27828v1 (CC BY 4.0), distributed by arXiv under the Creative Commons Attribution 4.0 International licence, http://creativecommons.org/licenses/by/4.0/, as recorded in the arXiv licence metadata for this exact version and shown on the abstract page https://arxiv.org/abs/2603.27828v1. Full licence notice: \texttt{licenses/arxiv-2603.27828-CC-BY-4.0.txt}.\par\smallskip

\noindent\textit{Editorial scope.} Counted unit: the article's proposition proposition (source0.tex lines 942--945). The article's complete original proof of it is delivered with it (source0.tex lines 946--1021, one \texttt{proof} environment of 76 lines). No mathematics is written, repaired or re-derived: the statement and the proof are the article's own lines, in the article's own notation, and no retained line is substituted or re-flowed.  Retained uncounted as the source-local context the proof needs: lines 562--568.  Nothing was omitted from any retained span, so the package carries no omission record.  No retained line contains a citation, so the wrapper carries no bibliography and no bibliography entry.  No retained line contains a figure environment, an image-inclusion command or drawing code, so no image is delivered and the package declares no asset dependency.  Editorial formatting only, in addition: an AMS \texttt{amsart} preamble that restates the article's own declarations and macros used by the retained lines, namely the environments \texttt{proposition} on separate counters, together with the standard packages those lines need.  The retained environments print in this document as Proposition 1 in order of appearance, whereas the article prints the same items with the numbering of its own full document.  The complete native source of the article as distributed in the arXiv e-print for this exact version is bundled unchanged as \texttt{sources/197385/source0.tex}.
\begin{equation}\label{eq:MPE}
    \begin{aligned}
        M^{n} =& \Delta x \Delta v \sum_{i,\bm{j}} f_{i,\bm{j}}^{n}, ~~ &\bm{P}^{n} =& \Delta x \Delta v \sum_{i,\bm{j}} \bm{v}_{\bm{j}} f_{i,\bm{j}}^{n}, \\
        \mathcal{E}_{K}^{n} =& \Delta x \Delta v \sum_{i,\bm{j}} \dfrac{|\bm{v}_{\bm{j}}|^{2}}{2} f_{i,\bm{j}}^{n}, ~~ &\mathcal{E}_{P}^{n} =& \Delta x \sum_{i} \dfrac{\lambda_{D}^{2}}{2} |E_{i}^{n}|^{2} , \\
        \mathcal{E}^{n} =& \mathcal{E}_{K}^{n} + \mathcal{E}_{P}^{n}, ~~ &S^{n} =& - \Delta x \Delta v \sum_{i,\bm{j}} f_{i,\bm{j}}^{n} \log f_{i,\bm{j}}^{n} .
    \end{aligned}
\end{equation}

\begin{proposition}
    This discretization of the Vlasov-Poisson equation preserves the discrete mass and momentum, while and kinetic energy s not conserved.
    % and satisfies a discrete entropy inequality under the CFL condition $\frac{\Delta t}{\Delta x}\max |v_x|\le 1$.
\end{proposition}

\begin{proof}
The discrete mass, momentum and energy are defined in Eq. \eqref{eq:MPE},
% \begin{equation}
%     \begin{aligned}
%         M^{n} =& \Delta x \Delta v \sum_{i,\bm{j}} f_{i,\bm{j}}^{n}, ~~ &\bm{P}^{n} =& \Delta x \Delta v \sum_{i,\bm{j}} \bm{v}_{\bm{j}} f_{i,\bm{j}}^{n}, \\
%         \mathcal{E}_{K}^{n} =& \Delta x \Delta v \sum_{i,\bm{j}} \dfrac{\bm{v}_{\bm{j}}^{2}}{2} f_{i,\bm{j}}^{n}, ~~ &\mathcal{E}_{P}^{n} =& \Delta x \sum_{i} \dfrac{\lambda_{D}^{2}}{2} |E_{i}^{n}|^{2} , \\
%         \mathcal{E}^{n} =& \mathcal{E}_{K}^{n} + \mathcal{E}_{P}^{n}, ~~ &S^{n} =& - \Delta x \Delta v \sum_{i,\bm{j}} f_{i,\bm{j}}^{n} \log f_{i,\bm{j}}^{n} .
%     \end{aligned}
% \end{equation}


\paragraph{Mass conservation}

\begin{equation}
    \dfrac{M^{n+1}-M^n}{\Delta t} = - \Delta x \Delta v \sum_{\bm{j}} v_{j_{x}} \sum_{i} D_{x}^{up} f_{i,\bm{j}}^{n} - \Delta x \Delta v \sum_{i} E_{i}^{n} \sum_{\bm{j}} D_{v}^{cen} f_{i,\bm{j}}^{n} = 0 ,
\end{equation}
where the first part vanishes due to telescoping cancellation, and the second part vanishes due to zero boundary condition.
% $f_{i,\bm{j}}^{n} = 0$ at boundary.


\paragraph{Momentum conservation}

\begin{equation}
\begin{aligned}
    \dfrac{\bm P^{n+1}-\bm P^n}{\Delta t} &= - \Delta x \Delta v \sum_{\bm{j}} v_{\bm{j}} v_{j_{x}} \sum_{i} D_{x}^{up} f_{i,\bm{j}}^{n} - \Delta x \Delta v \sum_{i} E_{i}^{n} \sum_{\bm{j}} v_{\bm{j}} D_{v}^{cen} f_{i,\bm{j}}^{n} \\
    &= \Delta x \sum_{i} E_{i}^{n} \rho_{i}^{n} = 0 ,
\end{aligned}
\end{equation}
where $\rho_{i}^{n} = \Delta v \sum_{\bm{j}} f_{i,\bm{j}}^{n}$.
The first term vanishes due to telescoping cancellation, and the second equation vanishes when we use FFT or central difference scheme on the Poisson equation.


\paragraph{Energy change}

\begin{equation}
    \dfrac{\mathcal E_K^{n+1}-\mathcal E_K^n}{\Delta t} = \Delta x \sum_{i} E_{i}^{n} J_{i}^{n},
\end{equation}
where $J_{i}^{n} = \Delta v \sum_{\bm{j}} \bm{v}_{\bm{j}} f_{i,\bm{j}}^{n}$, and it is similar in the momentum derivation.

The potential energy change is 
\begin{equation}
    \dfrac{\mathcal{E}_{P}^{n+1}-\mathcal{E}_{P}^n}{\Delta t} = \dfrac{\lambda_{D}^{2} \Delta x}{2 \Delta t} \sum_{i} (E_{i}^{n+1} + E_{i}^{n}) (E_{i}^{n+1} - E_{i}^{n}) ,
\end{equation}
and $E_{i}^{n+1} - E_{i}^{n} = - \lambda_{D}^{-2} J_{i}^{n} \Delta t + \cO(\Delta x \Delta t)$,
so that 
\begin{equation}
\begin{aligned}
    \dfrac{\mathcal{E}_{P}^{n+1}-\mathcal{E}_{P}^{n}}{\Delta t} =& \dfrac{\lambda_{D}^{2} \Delta x}{2 \Delta t} \sum_{i} [2E_{i}^{n} - \lambda_{D}^{-2} J_{i}^{n} \Delta t + \cO(\Delta x \Delta t)] [- \lambda_{D}^{-2} J_{i}^{n} \Delta t + \cO(\Delta x \Delta t)] \\
    =& - \Delta x \sum_{i} E_{i}^{n} J_{i}^{n} + \sum_{i} \dfrac{\lambda_{D}^{-2} \Delta t \Delta x}{2} |J_{i}^{n}|^{2} + \cO(\Delta x) + \cO(\Delta x \Delta t) .
\end{aligned}
\end{equation}

Therefore, there is a first-order time error in the total energy $\frac{\mathcal{E}_{P}^{n+1}-\mathcal{E}_{P}^{n}}{\Delta t} = \cO(\Delta t)$.
% because in the central-difference Poisson
% \begin{equation*}
% \begin{aligned}
%     (\rho_{i}^{n+1} - \rho_{i}^{n}) / \Delta t = - (J_{i}^{n} - J_{i-1}^{n}) / \Delta x + \cO(\Delta x) = - D_{x} J_{i}^{n} + \cO(\Delta x), \\
%     D_{x} (E_{i}^{n+1} - E_{i}^{n}) = \lambda_{D}^{-2} (\rho_{i}^{n+1} + \rho_{i-1}^{n+1} - \rho_{i}^{n} - \rho_{i-1}^{n})/2
% \end{aligned}
% \end{equation*}



% \paragraph{Entropy production.}
% Under the CFL condition $\frac{\Delta t}{\Delta x}|v_{j_x}|\le 1$, the update for each fixed velocity index $\bm j$ can be written as a convex combination
% \begin{equation*}
%     f_{i,\bm j}^{n+1} = \alpha_{i-1} f_{i-1,\bm j}^n + \alpha_i f_{i,\bm j}^n + \alpha_{i+1} f_{i+1,\bm j}^n, \quad \alpha_k\ge 0,\; \sum_k \alpha_k=1.
% \end{equation*}
% For any concave function $\eta$, Jensen's inequality implies
% \begin{equation*}
%     \eta(f_{i,\bm j}^{n+1}) \geq \alpha_{i-1}\eta(f_{i-1,\bm j}^n) + \alpha_i\eta(f_{i,\bm j}^n) + \alpha_{i+1}\eta(f_{i+1,\bm j}^n).
% \end{equation*}
% Summing over $i$ and $\bm j$ and choosing $\eta(f)=-f\log f$ gives the discrete entropy production.
% \end{proof}

\end{proof}

\end{document}
